\begin{proof}
\textbf{Step 1. Degenerate or edge cases.}  
If \(n=1\) then both sides of  
\[
n\sum_{i=1}^{n}a_{i}b_{i}=a_{1}b_{1}
\qquad\text{and}\qquad
\Bigl(\sum_{i=1}^{n}a_{i}\Bigr)\Bigl(\sum_{i=1}^{n}b_{i}\Bigr)=a_{1}b_{1}
\]
coincide, so the inequality holds with equality.  

If the \(a\)-sequence is constant, say \(a_{i}=a\) for all \(i\), then
\[
n\sum_{i=1}^{n}a_{i}b_{i}=na\sum_{i=1}^{n}b_{i}
   =\Bigl(\sum_{i=1}^{n}a_{i}\Bigr)\Bigl(\sum_{i=1}^{n}b_{i}\Bigr),
\]
and similarly when the \(b\)-sequence is constant.  Hence the statement is true in all degenerate situations.

\textbf{Step 2. Centered variables and means.}  
Define the arithmetic means
\[
\bar a:=\frac{1}{n}\sum_{i=1}^{n}a_{i},\qquad
\bar b:=\frac{1}{n}\sum_{i=1}^{n}b_{i},
\]
and the centred sequences
\[
\tilde a_{i}=a_{i}-\bar a,\qquad
\tilde b_{i}=b_{i}-\bar b\qquad(i=1,\dots ,n).
\]
Subtracting a constant does not affect the ordering, so
\[
\tilde a_{1}\le\tilde a_{2}\le\cdots\le\tilde a_{n},
\qquad
\tilde b_{1}\le\tilde b_{2}\le\cdots\le\tilde b_{n}. \tag{2}
\]
Moreover,
\[
\sum_{i=1}^{n}\tilde a_{i}=0,\qquad
\sum_{i=1}^{n}\tilde b_{i}=0. \tag{3}
\]

\textbf{Step 3. Non‑negativity of the covariance (Rearrangement Inequality).}  
For two sequences sorted in the same non‑decreasing order, the Rearrangement
Inequality states that for any permutation \(\sigma\) of \(\{1,\dots ,n\}\)
\[
\sum_{i=1}^{n}x_{i}y_{\sigma(i)}\le\sum_{i=1}^{n}x_{i}y_{i},
\]
with equality only when the sequences are linearly dependent after centring.  
Apply this to \((\tilde a_{i})\) and \((\tilde b_{i})\).  Because of (2) they are similarly ordered; for the reverse permutation
\(\rho(i)=n+1-i\) we obtain
\[
\sum_{i=1}^{n}\tilde a_{i}\tilde b_{\rho(i)}\le\sum_{i=1}^{n}\tilde a_{i}\tilde b_{i}. \tag{4}
\]
Using (3),
\[
\sum_{i=1}^{n}\tilde a_{i}\tilde b_{\rho(i)}
   =\sum_{i=1}^{n}\tilde a_{i}\bigl(\tilde b_{\,n+1-i}\bigr)
   =-\sum_{i=1}^{n}\tilde a_{i}\tilde b_{i},
\]
so (4) becomes
\[
-\sum_{i=1}^{n}\tilde a_{i}\tilde b_{i}\le\sum_{i=1}^{n}\tilde a_{i}\tilde b_{i},
\]
hence
\[
\boxed{\displaystyle\sum_{i=1}^{n}\tilde a_{i}\tilde b_{i}\ge0}. \tag{5}
\]
Thus the covariance of two similarly ordered centred sequences is non‑negative.

\textbf{Step 4. Algebraic manipulation to Chebyshev’s form.}  
Expand each product in (5):
\[
\tilde a_{i}\tilde b_{i}
   =(a_{i}-\bar a)(b_{i}-\bar b)
   =a_{i}b_{i}-a_{i}\bar b-\bar a b_{i}+\bar a\bar b.
\]
Summing over \(i\) and using linearity,
\[
\sum_{i=1}^{n}\tilde a_{i}\tilde b_{i}
   =\sum_{i=1}^{n}a_{i}b_{i}
    -\bar b\sum_{i=1}^{n}a_{i}
    -\bar a\sum_{i=1}^{n}b_{i}
    +n\bar a\bar b.
\]
Since \(\sum_{i=1}^{n}a_{i}=n\bar a\) and \(\sum_{i=1}^{n}b_{i}=n\bar b\),
\[
\sum_{i=1}^{n}\tilde a_{i}\tilde b_{i}
   =\sum_{i=1}^{n}a_{i}b_{i}
    -\frac{\bigl(\sum_{i=1}^{n}a_{i}\bigr)\bigl(\sum_{i=1}^{n}b_{i}\bigr)}{n}.
\]
Insert this expression into (5) to obtain
\[
\sum_{i=1}^{n}a_{i}b_{i}
   \ge\frac{\bigl(\sum_{i=1}^{n}a_{i}\bigr)\bigl(\sum_{i=1}^{n}b_{i}\bigr)}{n}.
\]
Multiplying by \(n\) yields Chebyshev’s sum inequality
\[
n\sum_{i=1}^{n}a_{i}b_{i}
   \ge\Bigl(\sum_{i=1}^{n}a_{i}\Bigr)\Bigl(\sum_{i=1}^{n}b_{i}\Bigr),
\]
which is precisely the desired statement \((1)\).

\textbf{Step 5. Equality conditions.}  
Equality in (5) occurs exactly when equality holds in the Rearrangement
Inequality for the centred sequences.  The Rearrangement Inequality is strict
unless there exists a constant \(\lambda\) such that
\[
\tilde a_{i}= \lambda\,\tilde b_{i}\qquad\text{for every }i. \tag{6}
\]
Re‑expressing (6) in the original variables gives
\[
a_{i}= \lambda b_{i}+\mu,\qquad
\mu:=\bar a-\lambda\bar b, \tag{7}
\]
so equality holds in one of the following mutually exclusive situations:
\begin{enumerate}
\item All \(a_{i}\) are equal (\(\tilde a_{i}=0\) for all \(i\));
\item All \(b_{i}\) are equal (\(\tilde b_{i}=0\) for all \(i\));
\item There exist constants \(\lambda\ge0\) and \(\mu\) such that
      \(a_{i}= \lambda b_{i}+\mu\) for every \(i\) (an affine relation preserving the order).
\end{enumerate}
These are exactly the equality cases for Chebyshev’s sum inequality.

\textbf{Step 6. Conclusion.}  
All degenerate cases have been handled, the centred covariance has been shown
to be non‑negative via the Rearrangement Inequality, the algebraic passage to
Chebyshev’s form is rigorous, and the equality conditions are completely
characterised.  Hence for any two non‑decreasing real sequences
\((a_i)_{i=1}^{n}\) and \((b_i)_{i=1}^{n}\) we have
\[
n\sum_{i=1}^{n}a_{i}b_{i}\;\ge\;
\Bigl(\sum_{i=1}^{n}a_{i}\Bigr)\Bigl(\sum_{i=1}^{n}b_{i}\Bigr),
\]
with equality precisely in the cases listed in Step 5. \(\square\)
\end{proof}